\scanpage{16}
il est immédiat qu'il l'est pour $(X,\Phi)$: si $(X'_j)_{j\in J}$ est la partition cherchée de $X'$, celle de $X$ est formée des $X'_j$ et des $U_i$. Ceci nous fournit un principe de dém. récurrente, si on suppose que $F$ satisfait la condition e)

\deleted{(Dim) Toute suite \inserted{\unsure{finie}} \deleted{décroissante}\inserted{$X_1\supset X_2\supset\cdots\supset X_n$} de parties fermées $\in F$ de $X$, telles que $\forall i$, $X_{i+1}$ soit non dense $X_i$, on a $X_n=\varnothing$ pour $n$ assez grand.}
\margin{\unsure{résulte en particulier de l'existence d'une dimension \unsure{finie}}}

\uline{Cor 1} Si $F$, \deleted{en plus des conditions}\inserted{satisfait les} a) b) c) \inserted{d) e)}\deleted{, satisfait la condition (Dim)}, alors $F$ est réunion filtrante des sous-ens. \inserted{de la forme} $F_{\mathcal P}$ (associés aux décompositions cellulaires \deleted{\unsure{compatibles}} $F$-admissibles de $X$).

\uline{Cor. 2} L'ens \inserted{$\unsure{\mathcal E|\mathcal D}$} des décompositions cellulaires $F$-admissibles, pour la relation d'ordre du raffinement, admet des bornes sup. pour les \inserted{sous} ens. finis de \deleted{\unsure{même}}\inserted{$\unsure{\mathcal E|\mathcal D}$}.
\mspage{17}
